% discussion.tex

\section{Discussion}

\subsection{Why Parametric EQ Improves Matching}

A subtractive synthesizer can only \emph{remove} spectral energy: the low-pass filter attenuates frequencies above its cutoff, but no component in the basic signal chain can add energy to a frequency band that the oscillators did not produce.
When the target sound contains a prominent presence peak in the 1--5\,kHz range---as our violin target does---the optimizer has no mechanism to reproduce it with oscillator mix and filter parameters alone.

Adding a two-band parametric equalizer resolves this asymmetry.
With just four extra parameters (center frequency and gain for each band), the EQ can \emph{boost} specific frequency regions after filtering.
In our experiments, CMA-ES converged on a +1\,dB boost near 5.8\,kHz, yielding an 8\% reduction in matching loss compared to the EQ-free baseline.
This result is notable because it came from the smallest parameter extension we tested (4 parameters), yet produced the largest improvement---confirming that \emph{targeted, physically motivated} additions to the signal chain outperform brute-force parameter expansion.

\subsection{Why More Parameters Hurt}

Expanding the synthesizer from 24 to 29 unconstrained parameters degraded performance.
CMA-ES exploited the additional degrees of freedom in unintended ways: the detune parameter drifted to $+24$ semitones (two octaves above the target fundamental), and the reverb mix climbed to 43\%, producing a washed-out output that happened to minimize the spectral loss without genuinely matching the target timbre.

This behavior illustrates a well-known tension in black-box optimization: more free parameters enlarge the feasible region, giving the optimizer room to find loss-minimizing solutions that do not correspond to perceptually meaningful patches.
The 29-parameter configuration effectively ``cheated'' the loss function rather than matching the sound.
The lesson is that \emph{constrained, physically meaningful parameters} are preferable to additional unconstrained ones, even when the latter could in principle increase expressiveness.

\subsection{The Harmonic Structure Gap}

Spectral analysis of our violin target reveals that the third harmonic ($H_3$) is stronger than the second ($H_2$).
This ordering is physically characteristic of bowed strings, where the bow's stick-slip excitation preferentially energizes odd-order modes.
However, it is \emph{impossible} to reproduce $H_3 > H_2$ with standard subtractive waveforms:
\begin{itemize}
    \item A sawtooth wave has all harmonics with amplitudes $\propto 1/n$, so $H_2 > H_3$ always.
    \item A square wave contains only odd harmonics ($H_1, H_3, H_5, \ldots$) with no $H_2$ at all, making the overall spectrum too hollow.
    \item No linear combination of saw and square can produce $H_3 > H_2 > 0$.
\end{itemize}

This harmonic ordering represents a hard floor for our subtractive architecture.
We tested three extensions designed to break through it---Chebyshev waveshaping (Test~A), FM synthesis (Test~H), and additive harmonic multipliers---but none improved the matching loss within 10K evaluations.
We attribute this to the loss function rather than the synthesis method: the multi-resolution STFT loss distributes error across all frequency bins and time frames, so a localized harmonic mismatch at $H_2$ and $H_3$ contributes only a small fraction of the total gradient signal.
A loss term that explicitly penalizes per-harmonic amplitude ratios may be needed to guide the optimizer toward the correct harmonic structure.

\subsection{Limitations}

Our study has several limitations that scope the generalizability of the results.

\paragraph{Single target sound.}
All experiments optimize against a single violin recording.
We have not evaluated whether the pipeline, loss weights, or parameter bounds transfer to other instruments (e.g., brass, woodwinds, plucked strings) or to synthetic targets where the ground-truth patch is known.

\paragraph{Subtractive synthesis only.}
The synthesizer architecture is limited to oscillator mixing, low-pass filtering, ADSR envelopes, and simple effects.
Timbres that require frequency modulation, wavetable interpolation, or physical modeling are out of scope for the current signal chain.

\paragraph{No real-time capability.}
Each CMA-ES evaluation renders a full note and computes a multi-resolution STFT loss, taking approximately 5--10\,ms per evaluation on CPU.
At 50K--500K evaluations per run, total optimization time ranges from minutes to hours, precluding interactive or real-time use.

\paragraph{Loss plateau.}
Beyond approximately 50K evaluations, we observe diminishing returns: the matching loss plateaus regardless of population size, step-size adaptation, or multi-start restarts.
This plateau suggests an architectural limit of the subtractive synthesizer rather than insufficient search---the optimizer has found the best patch the architecture can express.

\subsection{Future Work}

Several directions could address the limitations identified above.

\paragraph{Learned encoder.}
Replacing CMA-ES with a trained neural encoder (supervised, reinforcement learning, or flow matching) could amortize the optimization cost across many target sounds, enabling single-forward-pass inference.
Prior work on differentiable synthesizers~\cite{engel2020ddsp, masuda2023improving} demonstrates that encoder--decoder architectures can achieve competitive quality with orders-of-magnitude faster inference.

\paragraph{Richer synthesis architectures.}
Extending the signal chain with wavetable oscillators, FM operators, or Chebyshev waveshaping---paired with appropriate loss guidance---could break through the harmonic structure floor.
The key insight from our hypothesis tests is that adding synthesis capacity alone is insufficient; the loss function must also be able to ``see'' the harmonic mismatch.

\paragraph{Perceptual loss functions.}
Replacing or augmenting the multi-resolution STFT loss with perceptual metrics such as CDPAM~\cite{manocha2021cdpam} or learned temporal fine structure (TFS) losses could improve timbral fidelity.
Our TFS envelope loss experiment (Test~C) showed marginal improvement, but more sophisticated perceptual losses that operate on learned audio representations may better capture the nuances that spectral losses miss.

\paragraph{Real-time inference.}
Distilling the CMA-ES search into a lightweight neural network---or using gradient-based optimization on a fully differentiable synthesizer---could enable real-time parameter estimation suitable for live performance or interactive sound design tools.
